% -*- coding: utf-8 -*-

\chapter{实验与结果分析} \label{chpt:experiment}

\section{实验环境与数据集}

\subsection{实验环境}

% 表格形式列出实验环境：
% | 项目 | 配置 |
% |------|------|
% | 操作系统 | Windows 11 |
% | 编程语言 | Python 3.x |
% | 深度学习框架 | PyTorch 2.7.1 + CUDA 11.8 |
% | GPU | （填写实际GPU型号） |
% | 其他依赖 | torchvision 0.22.1, numpy, matplotlib |

\subsection{数据集}

% MNIST手写数字数据集：
% - 训练集：60,000张 28×28 灰度图像
% - 测试集：10,000张 28×28 灰度图像
% - 类别数：10类（数字0-9）
% - 归一化：均值=0.1307，标准差=0.3081
%
% 选择MNIST数据集的原因：
% 1. 经典的图像分类基准数据集
% 2. 计算效率高，便于大量实验对比
% 3. 对抗鲁棒性研究中广泛使用，便于横向比较

\subsection{基础模型}

% LeNet-5模型结构说明（可用表格或图示）
% 标准训练的基线性能：清洁样本准确率 99.0%


\section{评估指标与实验设置}

\subsection{评估指标}

% 1. 清洁准确率（Clean Accuracy）
%    在未施加任何攻击的原始测试样本上的分类准确率
%    衡量防御训练对模型正常性能的影响

% 2. 对抗准确率（Adversarial Accuracy）
%    在施加特定攻击后的对抗样本上的分类准确率
%    衡量模型对该攻击的防御能力

% 3. 攻击成功率（Attack Success Rate）
%    ASR = 1 - 对抗准确率
%    从攻击者视角衡量攻击效果

\subsection{攻击类型设置}

% 列出所有12种评估攻击及其参数：
%
% 梯度类攻击：
% - FGSM：ε=8/255
% - PGD (L∞)：ε=8/255, step_size=2/255, iterations=10
% - C&W (L∞)：ε=8/255, confidence=50
%
% 遮蔽类攻击：
% - 固定显著性遮蔽（k=3, 9, 15）
% - 自适应显著性遮蔽（N=3,R=2 / N=5,R=3 / N=10,R=3）
% - 固定IG遮蔽（对比基线）
% - 自适应IG遮蔽（对比基线）

\subsection{防御策略设置}

% 列出所有10种防御策略：
% 1. Standard：标准训练（无防御基线）
% 2. PGD-AT：PGD对抗性训练
% 3. FGSM-AT：FGSM对抗性训练
% 4. Occlusion-AT (IG)：基于梯度积分的遮蔽对抗性训练
% 5. Saliency-Occlusion-AT：基于显著性图的固定遮蔽对抗性训练
% 6. Adaptive-Saliency-AT：基于显著性图的自适应遮蔽对抗性训练
% 7. Adaptive-Saliency-AT (N=5,R=3)：自适应显著性训练（增强参数）
% 8. Adaptive-Occlusion-AT (IG)：基于梯度积分的自适应遮蔽训练
% 9. Mix-AT (Occlusion+PGD)：混合对抗性训练（基础版）
% 10. Adaptive-Mix-AT：自适应混合对抗性训练


\section{遮蔽攻击效果评估}

\subsection{固定遮蔽攻击与自适应遮蔽攻击的对比}

% 在标准模型上对比两种攻击方式的效果：
%
% | 攻击方式 | 参数 | 标准模型准确率 |
% |---------|------|--------------|
% | 无攻击 | - | 99.0% |
% | 固定遮蔽 | k=3 | 86.82% |
% | 固定遮蔽 | k=9 | 71.19% |
% | 固定遮蔽 | k=15 | 59.78% |
% | 自适应遮蔽 | N=3,R=2 | 49.45% |
% | 自适应遮蔽 | N=5,R=3 | 34.33% |
% | 自适应遮蔽 | N=10,R=3 | 15.95% |
%
% 分析：
% 1. 自适应遮蔽攻击效果显著强于固定遮蔽
% 2. 随参数增大，攻击效果增强但边际递减
% 3. N=5,R=3 与 N=10,R=3 差距说明增大区域数仍有明显效果

\subsection{遮蔽攻击的可视化分析}

% 插入对抗样本可视化图（数字0-9）：
% 图1：固定遮蔽攻击效果（不同k值）
% 图2：自适应遮蔽攻击效果（不同N,R值）
% 图3：不同遮蔽颜色的效果对比（黑色、灰色、白色）
%
% 分析可视化结果：
% - 遮蔽位置集中在笔画关键区域（转折点、交叉点等）
% - 自适应攻击的遮蔽更精确、面积更小
% - 黑色遮蔽效果最佳（与MNIST数据分布一致）

\subsection{与梯度类攻击的对比}

% 对比PGD、FGSM、C&W与遮蔽攻击对标准模型的攻击效果：
% - PGD：33.34%
% - FGSM：79.46%
% - C&W：33.37%
% - 自适应遮蔽(N=5,R=3)：34.33%
%
% 说明自适应遮蔽攻击效果与PGD攻击相当，
% 但扰动形式完全不同（局部遮蔽 vs 全局噪声）


\section{对抗性训练防御效果对比}

\subsection{综合对比实验结果}

% 核心实验结果表（10种防御 × 12种攻击的完整矩阵）
% 使用 comprehensive_multi_model_comparison.csv 的数据
%
% 建议使用横向大表或分成多个子表展示：
%
% 表1：各防御策略在梯度类攻击下的准确率
% | 防御策略 | Clean | FGSM | PGD | C&W |
% |---------|-------|------|-----|-----|
% | Standard | 99.0% | 79.46% | 33.34% | 33.37% |
% | PGD-AT | 99.3% | 96.1% | 94.95% | 95.05% |
% | FGSM-AT | 99.2% | 95.68% | 90.64% | 90.83% |
% | Occlusion-AT(IG) | 95.75% | 3.24% | 0% | 0% |
% | Adaptive-Saliency-AT(N=5,R=3) | 98.73% | 41.56% | 4.78% | 5.73% |
% | Mix-AT(OCC+PGD) | 98.53% | 93.5% | 91.7% | 91.7% |
% | Adaptive-Mix-AT | 98.9% | 94.01% | 91.48% | 91.4% |
%
% 表2：各防御策略在遮蔽类攻击下的准确率
% | 防御策略 | Fixed(k=9) | Adaptive(N=5,R=3) | Adaptive(N=10,R=3) |
% |---------|-----------|-------------------|-------------------|
% | Standard | 71.19% | 34.33% | 15.95% |
% | PGD-AT | 70.11% | 30.12% | 11.63% |
% | Occlusion-AT(IG) | 93.78% | 78.91% | 70.86% |
% | Adaptive-Saliency-AT(N=5,R=3) | 98.43% | 94.30% | 88.74% |
% | Mix-AT(Standard) | 95.48% | 80.65% | 54.7% |
% | Adaptive-Mix-AT | 94.03% | 76.06% | 50.66% |

\subsection{关键发现分析}

% 发现1：攻击类型特异性
% PGD-AT对梯度攻击防御优秀（94.95%），但对遮蔽攻击几乎无效（30.12%）
% 遮蔽-AT对遮蔽攻击防御优秀（94.30%），但对梯度攻击完全崩溃（0%）
% → 说明不同攻击类型需要不同的防御策略，验证了混合训练的必要性

% 发现2：自适应遮蔽训练效果最佳
% Adaptive-Saliency-AT(N=5,R=3)在遮蔽攻击下的防御效果：
% - 固定遮蔽(k=9)：98.43%（最高）
% - 自适应遮蔽(N=5,R=3)：94.30%（最高）
% - 自适应遮蔽(N=10,R=3)：88.74%（最高）
% → 自适应训练使模型学到了更鲁棒的特征表示

% 发现3：混合训练的均衡效果
% Adaptive-Mix-AT实现了最佳的综合防御：
% - 对PGD攻击：91.48%（接近PGD-AT的94.95%）
% - 对自适应遮蔽(N=5,R=3)：76.06%（显著高于PGD-AT的30.12%）
% - 清洁准确率：98.9%（几乎无损失）
% → 混合训练是多攻击场景下最实用的防御策略

% 发现4：清洁准确率影响
% 各防御策略的清洁准确率均在95%以上，
% 说明对抗性训练对正常分类性能的影响有限

\subsection{防御效果可视化}

% 插入多模型对比柱状图
% 图4：各防御策略在不同攻击下的准确率对比（分组柱状图）
% 图5：防御效果热力图（模型×攻击矩阵，颜色表示准确率）
% 图6：雷达图展示各防御策略的多维鲁棒性


\section{参数敏感性分析}

\subsection{遮蔽区域数量（top\_k / N）的影响}

% 在固定遮蔽攻击下，不同 top_k 值对攻击效果的影响：
% 图7：top_k vs 准确率 折线图
% - k=3: 86.82%
% - k=5: ~78%
% - k=9: 71.19%
% - k=15: 59.78%
% 分析：k值增大，攻击效果增强，但对防御模型的边际效果递减

% 在自适应遮蔽攻击下，不同 N 值的影响：
% 图8：N vs 准确率 折线图
% - N=3: 49.45%
% - N=5: 34.33%
% - N=10: 15.95%
% 分析：N的增大对攻击效果影响更显著

\subsection{遮蔽半径（R / kernel\_size）的影响}

% 图9：R vs 准确率 折线图
% - R=2: 49.15%
% - R=3: 34.33%
% - R=4: 23.91%
% 分析：半径增大导致遮蔽面积增大，攻击效果线性增强

\subsection{遮蔽颜色的影响}

% 对比三种遮蔽颜色（黑色0.0、灰色0.5、白色1.0）的效果
% 图10：遮蔽颜色 vs 准确率 柱状图
% 分析：黑色遮蔽效果最佳，白色和灰色效果相近
% 原因：MNIST数据集中背景为黑色，黑色遮蔽破坏前景特征更有效

\subsection{训练轮数的影响}

% 对比不同epoch下模型的防御效果：
% 图11：epoch vs 清洁准确率 + 对抗准确率 曲线
% 实验显示50 epoch后准确率趋于稳定，
% 继续训练到100/200 epoch无显著提升


\section{迁移攻击鲁棒性分析}

\subsection{白盒攻击与迁移攻击的对比}

% 实验设置：
% - 白盒攻击：使用目标模型自身生成对抗样本
% - 迁移攻击：使用标准模型生成对抗样本，在防御模型上测试
%
% 表3：白盒攻击 vs 迁移攻击对比
% | 攻击类型 | 白盒准确率 | 迁移准确率 |
% |---------|-----------|-----------|
% | FGSM | xx% | xx% |
% | PGD | xx% | xx% |
% | 固定遮蔽 | xx% | xx% |
% | 自适应遮蔽 | xx% | xx% |

\subsection{迁移攻击中的异常现象分析}

% 实验中发现的有趣现象：
% 对某些防御模型，迁移攻击的效果反而优于白盒攻击
% 例如：固定遮蔽攻击中，白盒攻击准确率94%，迁移攻击准确率91%
%
% 可能的解释：
% 1. 固定遮蔽的白盒攻击位置被防御模型"记住"并过度防御
% 2. 从标准模型生成的遮蔽位置与防御模型学到的防御模式不同
% 3. 攻击区域的重叠/分离特性影响攻击效果
%
% 可视化分析：
% 图12：白盒 vs 迁移攻击的遮蔽位置对比


\section{实验结果讨论}

% 1. 方法有效性总结
%    - 显著性图在遮蔽攻击中的有效性得到验证
%    - 自适应遮蔽攻击在所有攻击中效果最强
%    - 混合训练策略实现了最佳综合防御

% 2. 与现有方法的对比
%    - 对比PGD-AT等经典方法的局限性
%    - 说明本文方法的互补性和优越性

% 3. 方法的局限性
%    - 仅在MNIST数据集上验证，高分辨率数据集待验证
%    - 遮蔽攻击在某些场景下可能过于"明显"
%    - 混合训练需要更多计算资源


\section{本章小结}

% 本章在MNIST数据集上对所提出的基于显著性图的遮蔽攻击
% 及对抗性训练方法进行了全面的实验评估。
% 通过10种防御策略和12种攻击类型的系统对比，
% 验证了自适应遮蔽攻击的有效性和混合对抗性训练策略的优越性。
% 参数敏感性分析和迁移攻击实验进一步揭示了
% 遮蔽攻击参数对防御效果的影响规律和模型的迁移鲁棒性特征。
